\begin{helper}
Let \(H\) be a history cell, \(T\) a target event, \(U\) the terminal
coordinate set, and let
\[
G_{B,J,\beta,\tau}
\]
be a red branch/transcript cell family indexed by
\[
B\in\mathcal B,\qquad J\in\mathcal J,\qquad
\beta\in\mathcal T,\qquad \tau=(P_\tau,A_\tau,Z_\tau)\in\mathcal P,
\]
where \(\mathcal T\) is the finite branch/trace label type.
Assume \(0\le p\le 1/2\).  Assume the complete terminal product law over
\(H\): for every \(A\subseteq U\), the conditional probability, given \(H\),
of the complete assignment \(X_e=1\) exactly on \(A\) is
\[
p^{|A|}(1-p)^{|U|-|A|}.
\]
Assume also that whenever a sample lies in \(G_{B,J,\beta,\tau}\), the sets
\(P_\tau\) and \(A_\tau\) are disjoint subsets of \(U\), every coordinate in
\(P_\tau\) is present, and every coordinate in \(A_\tau\) is absent.

Then there is a real-valued cell weight \(w_{B,J,\beta,\tau}\), defined for
all \(B\in\mathcal B\), \(J\in\mathcal J\),
\(\beta\in\mathcal T\), and \(\tau\in\mathcal P\), such that for every
such index
\[
0\le w_{B,J,\beta,\tau},
\]
\[
w_{B,J,\beta,\tau}\le p^{|P_\tau|}(1-p)^{|A_\tau|},
\]
and the raw event-probability bound holds:
\[
\Pr(T\cap H\cap G_{B,J,\beta,\tau})
\le
\Pr(H)\,w_{B,J,\beta,\tau}.
\]
\end{helper}
\begin{proof}
Fix indices \(B,J,\beta,\tau\), with
\(\tau=(P_\tau,A_\tau,Z_\tau)\).  Set
\[
w_{B,J,\beta,\tau}=p^{|P_\tau|}(1-p)^{|A_\tau|}.
\]
Since \(0\le p\le 1/2\), also \(p\le1\) and \(0\le1-p\).  Hence this weight
is nonnegative and it satisfies the displayed cylinder-factor upper bound
with equality.

The cell-geometry hypothesis gives, for every sample in the cell,
\[
P_\tau,A_\tau\subseteq U,\qquad P_\tau\cap A_\tau=\emptyset,
\]
and the cell event implies that every coordinate of \(P_\tau\) is present
and every coordinate of \(A_\tau\) is absent.  Hence
\[
T\cap H\cap G_{B,J,\beta,\tau}
\subseteq
H\cap\{X_e=1\text{ for }e\in P_\tau,\,
      X_e=0\text{ for }e\in A_\tau\}.
\]
The right-hand event is a terminal cylinder: it is the disjoint union, over
all complete terminal assignments \(S\subseteq U\) with
\(P_\tau\subseteq S\) and \(S\cap A_\tau=\emptyset\), of the complete
assignment events \(X_e=1\) exactly for \(e\in S\).  The complete terminal
product law assigns such an \(S\) the conditional mass
\[
p^{|S|}(1-p)^{|U|-|S|}.
\]
Factoring out the forced present coordinates in \(P_\tau\) and the forced
absent coordinates in \(A_\tau\), the unmentioned terminal coordinates each
contribute \(p+(1-p)=1\).  Thus the conditional mass of the whole cylinder
is bounded by, and in the exact product case equals,
\[
p^{|P_\tau|}(1-p)^{|A_\tau|}.
\]

Apply the terminal cylinder charge
\(\noderef{FixedSetHistoryCellFixedCylinderCharge}\), using the complete
terminal product law for \(H\), to the disjoint present and absent sets
\(P_\tau,A_\tau\).  It gives the conditional cylinder bound
\[
\Pr(P_\tau\text{ present and }A_\tau\text{ absent}\mid H)
\le p^{|P_\tau|}(1-p)^{|A_\tau|}.
\]

It remains to pass carefully from this conditional estimate to the raw cell
mass estimate.  Write \(\Pr\) for
\(\noderef{TwoBiteEventProbability}\).  The sample weights are nonnegative
by \(\noderef{TwoBiteSampleWeightNonnegative}\), since \(0\le p\le1\).  Hence
event probability is monotone under containment: if \(E\subseteq F\), then
the finite sum defining \(\Pr(E)\) is obtained from the sum for \(\Pr(F)\)
by deleting nonnegative terms, so \(\Pr(E)\le\Pr(F)\).

First suppose \(\Pr(H)>0\).  By the definition of
\(\noderef{TwoBiteConditionalProbability}\), the conditional probability of
the cylinder given \(H\) is
\[
\frac{\Pr(H\cap\{P_\tau\text{ present},A_\tau\text{ absent}\})}{\Pr(H)}.
\]
The displayed conditional cylinder bound therefore gives
\[
\Pr(H\cap\{P_\tau\text{ present},A_\tau\text{ absent}\})
\le
\Pr(H)\,w_{B,J,\beta,\tau}
\]
after multiplying by the positive denominator \(\Pr(H)\).  By the containment
of \(T\cap H\cap G_{B,J,\beta,\tau}\) in this cylinder event and by
monotonicity,
\[
\Pr(T\cap H\cap G_{B,J,\beta,\tau})
\le \Pr(H)\,w_{B,J,\beta,\tau}.
\]

Now suppose \(\Pr(H)=0\).  The event \(T\cap H\cap G_{B,J,\beta,\tau}\) is a
subevent of \(H\).  Monotonicity from the nonnegative sample weights gives
\[
\Pr(T\cap H\cap G_{B,J,\beta,\tau})\le \Pr(H)=0.
\]
The left side is also nonnegative, again by nonnegative sample weights, so it
is \(0\).  Since \(w_{B,J,\beta,\tau}\ge0\), the right side
\(\Pr(H)w_{B,J,\beta,\tau}\) is also \(0\), giving the raw bound in the
zero-history-mass case.  This two-case multiplication step is the concrete
case split behind the general monotone conditional-intersection estimate in
\(\noderef{TwoBiteConditionalIntersectionBound}\).
\end{proof}
